\documentclass[11pt,a4paper]{article}

% Sprache und Schriften (mit XeLaTeX kompilieren)
\usepackage{fontspec}
\usepackage[ngerman]{babel}

\usepackage[margin=2.5cm]{geometry}
\usepackage{amsmath, amssymb}
\usepackage{graphicx}
\usepackage{enumitem}
\usepackage[hidelinks]{hyperref}

% Angaben zur Aufgabe
\newcommand{\kurs}{Name der Vorlesung}
\newcommand{\blatt}{Übungsblatt 1}
\newcommand{\student}{Vorname Nachname}
\newcommand{\dozent}{Name der Lehrperson}

% Nummerierte Frage und Umgebung für die Antwort
\newcounter{frage}
\newcommand{\frage}[1]{%
  \stepcounter{frage}%
  \section*{Frage \thefrage\ifx\relax#1\relax\else: #1\fi}}
\newenvironment{antwort}{\par\noindent\textbf{Antwort.}\ }{\hfill$\blacksquare$\par}

\begin{document}

\noindent
\begin{minipage}[t]{0.6\textwidth}
  \textbf{\kurs}\\
  \blatt
\end{minipage}%
\begin{minipage}[t]{0.4\textwidth}
  \raggedleft
  \student\\
  Lehrperson: \dozent\\
  \today
\end{minipage}
\par\medskip\hrule\medskip

\frage{Ableitungen}
Berechnen Sie $\dfrac{d}{dx}\left(x^2 \sin x\right)$.

\begin{antwort}
  Nach der Produktregel gilt
  \[
    \frac{d}{dx}\left(x^2 \sin x\right) = 2x \sin x + x^2 \cos x.
  \]
\end{antwort}

\frage{}
Nennen Sie zwei Eigenschaften des bestimmten Integrals.

\begin{antwort}
  \begin{enumerate}[label=(\alph*)]
    \item Linearität: $\int_a^b (\alpha f + \beta g) = \alpha \int_a^b f + \beta \int_a^b g$.
    \item Additivität: $\int_a^b f = \int_a^c f + \int_c^b f$.
  \end{enumerate}
\end{antwort}

\end{document}
